\documentclass{../latex/utmreport}
\newcommand{\ReportTitle}{Lab 5: The File System with Your Own Hands}
% Shared title-page and PDF metadata for every laboratory report.
\newcommand{\StudentName}{Gafenco Victor}
\newcommand{\StudentGroup}{FAF-241}
\newcommand{\CourseName}{Operating Systems}
\newcommand{\InstructorName}{Patricia Reițman}
\newcommand{\InstructorTitle}{Assistant lecturer}
\newcommand{\ReportCity}{Chișinău}
\newcommand{\ReportYear}{2026}
\newcommand{\CourseRepository}{https://github.com/yoda-dinmd/os-labs}

\selectlanguage{english}
\utmlogo{../latex/assets/UTM_Logo_English.jpeg}
\institution{Technical University of Moldova}
\faculty{Faculty of Computers, Informatics, and Microelectronics}
\department{Department of Software Engineering}
\studyprogramme{Software Engineering Study Programme}
\documenttype{Laboratory report}
\student{\StudentName}{\StudentGroup}
\supervisor{\InstructorName}{\InstructorTitle}
\renewcommand{\UTMCoordinatorLabel}{Instructor}
\cityyear{\ReportCity}{\ReportYear}
\reportsubtitle{\CourseName}
\hypersetup{pdfauthor={\StudentName},pdfsubject={\CourseName},
  pdftitle={\ReportTitle}}
\reporttitle{\ReportTitle}
\addto\captionsenglish{\renewcommand{\figurename}{Figure}\renewcommand{\tablename}{Table}\renewcommand{\bibname}{References}}
\renewcommand{\figurename}{Figure}
\renewcommand{\tablename}{Table}
\renewcommand{\bibname}{References}
\usepackage{float}
\usepackage{placeins}
\setlength{\emergencystretch}{2em}
\newcommand{\code}[1]{\texttt{\detokenize{#1}}}
% Keep screenshots beside their observations, with bounded height and width.
\newcommand{\screenshot}[4][\textwidth]{%
  \begin{figure}[H]
    \centering
    \includegraphics[width=#1,height=0.62\textheight,keepaspectratio]{screenshots/#2}
    \caption{#3}\label{#4}
  \end{figure}}
\newcommand{\reportreferences}[2]{%
  \clearpage\phantomsection\addcontentsline{toc}{chapter}{References}
  \begin{thebibliography}{9}
    \bibitem{handout} Technical University of Moldova, FCIM / FAF.
    \textit{#1}. Operating Systems laboratory handout, academic year 2026--2027.
    Supplied course document: \code{#2}.
  \end{thebibliography}}

\providecommand{\tightlist}{\setlength{\itemsep}{0pt}\setlength{\parskip}{0pt}}
\begin{document}
\makeutmreporttitlepage
\tableofcontents
\clearpage
\begin{utmabstract}[Abstract]
This laboratory investigates Linux file types, inode identity, links, permissions and storage. Experiments on openSUSE Tumbleweed establish the difference between names and file contents, demonstrate directory search permissions and sparse files, and map the root filesystem and live process information. The observed Btrfs directory link count is compared with the traditional Unix model.
\utmkeywords{Linux, Btrfs, inode, hard link, symbolic link, permissions, procfs}
\end{utmabstract}
\utmintroduction

This laboratory examines Linux file types, inode identity, hard and
symbolic links, permissions, disk layout and the process information
exposed by procfs. The experiments were performed on 12 October 2026 as
the ordinary user \texttt{victor} on openSUSE Tumbleweed. A fresh
directory, \texttt{/home/victor/os-lab5/run-018y4c78}, kept the file
experiments separate from existing work. Its filesystem is Btrfs with a
reported block size of 4096 bytes.

The following sections answer every Observe question in the handout.
Command excerpts show the measured results; the complete transcripts are
retained in \texttt{evidence/1-tree.txt} through
\texttt{evidence/5-stretch.txt}. Two optional experiments were
completed: recognizing file signatures and reading an unlinked file
through an existing descriptor.

\chapter{The tree and file types}\label{the-tree-and-file-types}

\section{Recognizing objects}\label{recognizing-objects}

The command \texttt{ls\ /} showed the usual top-level names, including
\texttt{bin}, \texttt{boot}, \texttt{dev}, \texttt{etc}, \texttt{home},
\texttt{proc}, \texttt{tmp}, \texttt{usr} and \texttt{var}. The
requested detailed listing was:

\noindent\begin{minipage}{\linewidth}
\begin{lstlisting}
$ ls -ld /etc /etc/passwd /bin /dev/null /tmp
lrwxrwxrwx.  1 root root    7 Mar  2  2026 /bin -> usr/bin
crw-rw-rw-.  1 root root 1, 3 Oct  6 10:15 /dev/null
drwxr-xr-x.  1 root root 3654 Oct 12 15:12 /etc
-rw-r--r--.  1 root root 2202 Oct  5 14:07 /etc/passwd
drwxrwxrwt. 73 root root 3900 Oct 12 16:10 /tmp
\end{lstlisting}
\end{minipage}\par\medskip

The first character identifies \texttt{/etc} and \texttt{/tmp} as
directories (\texttt{d}), \texttt{/etc/passwd} as a regular file
(\texttt{-}), \texttt{/bin} as a symbolic link (\texttt{l}), and
\texttt{/dev/null} as a character device (\texttt{c}). The trailing dot
after the mode is additional security-metadata notation, not a file type
or permission bit.

\texttt{/bin} stores the relative target \texttt{usr/bin}. Since the
link is in \texttt{/}, it resolves to \texttt{/usr/bin}.
\texttt{file\ /usr/bin/ls\ /etc/passwd\ /bin} identified an x86-64 ELF
executable, ASCII text, and a symbolic link respectively.

A character-device example is \texttt{/dev/null}, with major/minor
numbers 1 and 3. The block-device search
\texttt{find\ /dev\ -maxdepth\ 2\ -type\ b\ -ls} found
\texttt{/dev/nvme0n1} (259, 0) and its partitions, including
\texttt{/dev/nvme0n1p5} (259, 5). Their \texttt{b} mode prefix
distinguishes block devices from character devices.

\section{Comparing allocated space}\label{comparing-allocated-space}

\noindent\begin{minipage}{\linewidth}
\begin{lstlisting}
$ du -sh /var/log 2>/dev/null
320M    /var/log
$ du -sh /usr 2>/dev/null
13G     /usr
\end{lstlisting}
\end{minipage}\par\medskip

The visible allocation under \texttt{/usr} is approximately 42 times
that under \texttt{/var/log}: \texttt{13\ *\ 1024\ /\ 320\ =\ 41.6}. The
difference is about 12.7 GiB. These are rounded values and lower bounds
for the accessible directory contents: errors were suppressed as
instructed, and some contents were inaccessible to the ordinary user.
\texttt{du} measures allocated space, not the sum of logical file
lengths.

\chapter{Inodes and links}\label{inodes-and-links}

\section{Names and shared identity}\label{names-and-shared-identity}

After \texttt{echo\ hello\ \textgreater{}\ notes.txt},
\texttt{stat\ notes.txt} reported size 6 bytes, inode \textbf{468755},
and \textbf{Links: 1}. \texttt{Inode} identifies the inode within its
filesystem; \texttt{Links} counts hard links to it. The sixth byte is
the newline written by \texttt{echo}.

\noindent\begin{minipage}{\linewidth}
\begin{lstlisting}
$ ln notes.txt hard.txt
$ ls -li notes.txt hard.txt
468755 -rw-r--r--. 2 victor victor 6 Oct 12 16:10 hard.txt
468755 -rw-r--r--. 2 victor victor 6 Oct 12 16:10 notes.txt
$ rm notes.txt
$ cat hard.txt
hello
$ ls -li hard.txt
468755 -rw-r--r--. 1 victor victor 6 Oct 12 16:10 hard.txt
\end{lstlisting}
\end{minipage}\par\medskip

Both names refer to the same inode and the same data. Removing
\texttt{notes.txt} removes that directory entry and reduces the link
count from 2 to 1. It does not privilege the first name or erase the
contents while \texttt{hard.txt} still refers to them. After the last
name is removed, open references can still keep the file alive; the
optional experiment below demonstrates this.

\section{Symbolic links store a
pathname}\label{symbolic-links-store-a-pathname}

\noindent\begin{minipage}{\linewidth}
\begin{lstlisting}
$ ln -s hard.txt soft.txt
$ ls -li hard.txt soft.txt
468755 -rw-r--r--. 1 victor victor 6 Oct 12 16:10 hard.txt
468756 lrwxrwxrwx. 1 victor victor 8 Oct 12 16:10 soft.txt -> hard.txt
$ cat soft.txt
hello
$ rm hard.txt
$ cat soft.txt
cat: soft.txt: No such file or directory
$ ls -l soft.txt
lrwxrwxrwx. 1 victor victor 8 Oct 12 16:10 soft.txt -> hard.txt
\end{lstlisting}
\end{minipage}\par\medskip

The symbolic link has its own inode, 468756. Its length is exactly
\textbf{8 bytes}, the number of characters in \texttt{hard.txt}, with no
stored newline counted in its length. The pathname is resolved relative
to the directory containing the link. Removing \texttt{hard.txt} leaves
a dangling link: the link itself exists, but its target no longer
resolves. A symbolic link does not increase its target's hard-link
count.

\section{Directory link counts: prediction and
observation}\label{directory-link-counts-prediction-and-observation}

The handout's traditional prediction was 2 links for a new directory and
4 after adding two child directories. In that model, the parent's entry
for \texttt{d} and \texttt{d/.} account for two, and \texttt{d/x/..} and
\texttt{d/y/..} add two more.

The actual Btrfs result differs:

\noindent\begin{minipage}{\linewidth}
\begin{lstlisting}
$ mkdir d; ls -ld d
drwxr-xr-x. 1 victor victor 0 Oct 12 16:10 d
$ mkdir d/x d/y; ls -ld d
drwxr-xr-x. 1 victor victor 4 Oct 12 16:10 d
$ stat -f -c 'work filesystem: %T, block size %S' .
work filesystem: btrfs, block size 4096
\end{lstlisting}
\end{minipage}\par\medskip

The link-count column remains \textbf{1}. The final \texttt{4} is the
directory's size column, not its link count. Btrfs does not expose the
traditional \texttt{2\ +\ number\ of\ subdirectories} count here. Thus
the handout explains the traditional directory model, while this
laptop's measurement must be reported as 1 and 1.

\section{Logical length and allocated
storage}\label{logical-length-and-allocated-storage}

\noindent\begin{minipage}{\linewidth}
\begin{lstlisting}
$ echo x > one.txt; ls -l one.txt; du -h one.txt
-rw-r--r--. 1 victor victor 2 Oct 12 16:10 one.txt
4.0K    one.txt
$ truncate -s 1G big.img; ls -lh big.img; du -h big.img
-rw-r--r--. 1 victor victor 1.0G Oct 12 16:10 big.img
0       big.img
\end{lstlisting}
\end{minipage}\par\medskip

\texttt{one.txt} contains \texttt{x} and a newline but its reported
allocation occupies a 4096-byte unit. The sparse file has a logical
length of 1 GiB (1,073,741,824 bytes); \texttt{truncate} extends it with
holes rather than writing data blocks. Reading those holes yields zero
bytes. \texttt{du} reports zero allocated file-data space at this point;
filesystem metadata still exists. \texttt{rm\ big.img} removed the
experimental file.

\chapter{Permissions}\label{permissions}

\section{Modes and ownership}\label{modes-and-ownership}

The nine permission characters below omit the separate file-type
character.

{
\begin{longtable}[]{@{}lll@{}}
\hline
Octal & Permissions & Example use \\
\hline
\endhead
\hline
\endlastfoot
600 & \texttt{rw-\/-\/-\/-\/-\/-\/-} & A private configuration file \\
644 & \texttt{rw-r-\/-r-\/-} & A document others may read \\
755 & \texttt{rwxr-xr-x} & A public directory or executable \\
700 & \texttt{rwx-\/-\/-\/-\/-\/-} & A private directory \\
\end{longtable}
}

\noindent\begin{minipage}{\linewidth}
\begin{lstlisting}
$ echo secret > s.txt
$ chmod 600 s.txt; ls -l s.txt
-rw-------. 1 victor victor 7 Oct 12 16:10 s.txt
$ chmod 000 s.txt; cat s.txt
cat: s.txt: Permission denied
$ chmod u+r s.txt; cat s.txt
secret
\end{lstlisting}
\end{minipage}\par\medskip

An ordinary owner's access is checked against the owner permission bits.
Ownership permits changing the mode, but it does not itself grant read
access when those bits are zero. Restoring owner read permission allows
the file to be read again.

\section{The creation mask}\label{the-creation-mask}

\noindent\begin{minipage}{\linewidth}
\begin{lstlisting}
$ umask
0022
$ touch new.txt; mkdir newdir; ls -ld new.txt newdir
-rw-r--r--. 1 victor victor 0 Oct 12 16:10 new.txt
drwxr-xr-x. 1 victor victor 0 Oct 12 16:10 newdir
\end{lstlisting}
\end{minipage}\par\medskip

With no overriding default ACL, the creation mode is the requested mode
with the mask bits cleared. A normal file starts from \texttt{0666},
producing \texttt{0644} under \texttt{0022}; a directory starts from
\texttt{0777}, producing \texttt{0755}. The operation is a bit mask, not
general arithmetic subtraction. A new ordinary file does not acquire
execute bits merely because the mask allows them.

\section{Reading and searching a
directory}\label{reading-and-searching-a-directory}

Before each test, the prediction distinguished listing names from
looking up a known name. The file \texttt{box/f} was created with
\texttt{echo\ hi\ \textgreater{}\ box/f}.

\noindent\begin{minipage}{\linewidth}
\begin{lstlisting}
$ chmod u=r,go= box; ls box; cat box/f
f
cat: box/f: Permission denied
$ chmod u=x,go= box; ls box; cat box/f
ls: cannot open directory 'box': Permission denied
hi
$ chmod 755 box
\end{lstlisting}
\end{minipage}\par\medskip

Directory \textbf{read} permission permits enumerating entry names.
Directory \textbf{execute}, also called search permission, permits
traversing it and resolving a known entry. Consequently a read-only
directory can reveal \texttt{f} without allowing it to be opened; a
search-only directory allows \texttt{box/f} to be opened when its name
is known and the file's own mode permits reading, but prevents listing.
Directory write permission, normally combined with search, controls
creating and removing entries.

\section{Set-user-ID, the sticky bit and protected
files}\label{set-user-id-the-sticky-bit-and-protected-files}

\noindent\begin{minipage}{\linewidth}
\begin{lstlisting}
$ ls -l /usr/bin/passwd
-rwsr-xr-x. 1 root shadow 97408 Aug 10 09:57 /usr/bin/passwd
$ ls -ld /tmp
drwxrwxrwt. 77 root root 3980 Oct 12 16:10 /tmp
$ cat /etc/shadow
cat: /etc/shadow: Permission denied
$ ls -l /etc/shadow
-rw-r-----. 1 root shadow 977 Oct 5 14:07 /etc/shadow
\end{lstlisting}
\end{minipage}\par\medskip

The denied read was recorded in an earlier check as the same ordinary user.
The final run also confirmed that the file was not readable.

The \texttt{s} occupies the owner's execute position: owner execution
and the set-user-ID bit are both set. Under normal execution conditions,
\texttt{passwd} runs with its owner's effective user ID, root, allowing
controlled updates to protected password data after its own checks. The
user cannot directly read \texttt{/etc/shadow}; its mode grants access
only to root and the shadow group.

The \texttt{t} in \texttt{/tmp} means other-execute plus the sticky bit.
Even though many users can create entries in this directory, removal or
renaming of someone else's entries is restricted to the file owner,
directory owner or a suitably privileged process. It does not make file
contents private; file modes still govern that.

\chapter{Disks and the kernel as
files}\label{disks-and-the-kernel-as-files}

\section{Disk layout and root usage}\label{disk-layout-and-root-usage}

\texttt{lsblk\ -o\ NAME,TYPE,FSTYPE,SIZE,MOUNTPOINTS} established the
following layout. All partitions belong to the same 931.5 GiB NVMe disk;
repeated Btrfs mountpoints are not additional disks.

\noindent\begin{minipage}{\linewidth}
\begin{lstlisting}
nvme0n1                 931.5G  disk
|-- nvme0n1p1   vfat      200M  /boot/efi
|-- nvme0n1p2             16M   no mounted filesystem shown
|-- nvme0n1p3   ntfs    243.3G  not mounted
|-- nvme0n1p4   ntfs      846M  not mounted
|-- nvme0n1p5   btrfs   685.2G  /, /home, /var, ...
`-- nvme0n1p6   swap        2G  [SWAP]
\end{lstlisting}
\end{minipage}\par\medskip

The root mount is \texttt{/dev/nvme0n1p5}, filesystem type
\textbf{Btrfs}, with \texttt{rw} in its mount options. The selected root
subvolume is \texttt{/@/.snapshots/40/snapshot}, subvolume ID 303. The
\texttt{df\ -h} row was:

\noindent\begin{minipage}{\linewidth}
\begin{lstlisting}
Filesystem      Size  Used Avail Use% Mounted on
/dev/nvme0n1p5  686G   80G  605G  12% /
\end{lstlisting}
\end{minipage}\par\medskip

Thus root was \textbf{12\% full}, using approximately 80 GiB, with 605
GiB available. Human-readable rounding explains why the displayed sizes
need not add exactly.

\section{CPU, memory and the meaning of
self}\label{cpu-memory-and-the-meaning-of-self}

\noindent\begin{minipage}{\linewidth}
\begin{lstlisting}
$ grep -c processor /proc/cpuinfo
8
$ head -3 /proc/meminfo
MemTotal:       16212804 kB
MemFree:         1033936 kB
MemAvailable:    7479348 kB
$ head -5 /proc/self/status
Name:   head
Umask:  0022
State:  R (running)
Tgid:   842397
Ngid:   0
\end{lstlisting}
\end{minipage}\par\medskip

The processor count establishes \textbf{8 logical CPUs}, not necessarily
eight physical cores; simultaneous multithreading can expose multiple
logical CPUs per core. The kernel reports \textbf{16,212,804 KiB},
approximately \textbf{15.46 GiB}, of total usable RAM. This is distinct
from currently free or available memory and may be slightly below the
installed module capacity.

The process name is \textbf{head} because \texttt{/proc/self} resolves
for the process accessing the pathname. Here \texttt{head} opens and
reads its own status, rather than the shell's status. Similarly,
\texttt{ls\ -l\ /proc/self/fd} listed the descriptors of \texttt{ls}: 0,
1 and 2 referred to pipes in this run, and descriptor 3 referred to
\texttt{/proc/842398/fd}, opened while listing that directory. These are
live kernel views, not ordinary stored files.

\chapter{Optional experiments}\label{optional-experiments}

\section{File signatures rather than
suffixes}\label{file-signatures-rather-than-suffixes}

\noindent\begin{minipage}{\linewidth}
\begin{lstlisting}
$ gzip -c /etc/passwd > notes.txt
$ file notes.txt
notes.txt: gzip compressed data, was "passwd", ...
$ xxd notes.txt | head -1
00000000: 1f8b 0808 f584 c36a 0003 7061 7373 7764  .......j..passwd
\end{lstlisting}
\end{minipage}\par\medskip

The first two bytes, \textbf{1f 8b}, are the gzip signature. The next
byte, \texttt{08}, denotes the DEFLATE compression method. The
\texttt{.txt} suffix does not change the content: \texttt{file}
recognizes the header. The ellipsis shortens the descriptive timestamp
fields; the complete output is in the transcript.

\section{Removing a file that remains
open}\label{removing-a-file-that-remains-open}

\noindent\begin{minipage}{\linewidth}
\begin{lstlisting}
$ printf 'open file remains readable\n' > bigfile
$ exec 3< bigfile
$ rm bigfile
$ ls -l /proc/$$/fd
# Descriptor 3 points to:
# /home/victor/os-lab5/run-018y4c78/bigfile (deleted)
$ cat <&3
open file remains readable
$ exec 3<&-
\end{lstlisting}
\end{minipage}\par\medskip

The shell expands \texttt{\$\$} to its own process ID, so this time the
directory describes the shell's descriptors. The link ends in
\textbf{(deleted)} because the name has been unlinked, but descriptor 3
still holds the open file. Its contents remain readable. Closing
descriptor 3 releases that reference; with no links or other open
references remaining, the file's storage becomes reclaimable. On
filesystems with snapshots or shared extents, physical space can remain
referenced elsewhere. The subsequent descriptor listing contained only
0, 1 and 2.

\utmconclusions

The experiments separate names from inode identity, logical file length
from allocated storage, and file permissions from directory search
permissions. Hard links preserve access to the same inode after one name
is removed, while symbolic links can remain present but dangling. The
Btrfs directory-link result also shows why measured output must take
precedence over a generic filesystem example. Disk and procfs inspection
connected these abstractions to the laptop's root partition, logical
CPUs, memory and live processes.

\end{document}
